\documentclass[10pt,a4paper]{amsart}
\usepackage[margin=19mm]{geometry}
\usepackage{amsmath,amssymb,mathtools}
\usepackage[scr=rsfs,cal=euler]{mathalfa}
\usepackage{xfrac}
\usepackage[unicode,hidelinks,bookmarks=false]{hyperref}
\newcommand{\ie}{i.e.\ }
\newcommand{\cf}{cf.\ }
\newcommand{\resp}{respectively\ }
\newcommand{\Subsection}{\subsection}
\providecommand{\nobreakdash}{\nobreak\hbox{-}}
\setlength{\emergencystretch}{2em}
\multlinegap0pt
\usepackage{bm}
\usepackage{bbm}
\usepackage{dsfont}
\usepackage{xspace}
\usepackage{cancel}
\usepackage{mathtools}
\usepackage{amsthm}
\makeatletter
\@ifundefined{subsection}{\newtheorem{subsection}{Subsection}}{}
\@ifundefined{prop}{\newtheorem{prop}[subsection]{Proposition}}{}
\@ifundefined{thm}{\newtheorem{thm}[subsection]{Theorem}}{}
\makeatother
\newcommand{\Cs}{\mathsf{C}}
\newcommand{\Rs}{\mathsf{R}}
\renewcommand{\and}{\quad\text{and}\quad}
\title[Minimal $a$-numbers of Artin--Schreier covers of ordinary cu]{Minimal $a$-numbers of Artin--Schreier covers of ordinary curves}
\author{Bryden Cais, Douglas Ulmer}
\date{Source: arXiv:2604.20091v1 (CC BY 4.0)}
\begin{document}
\maketitle

\noindent\emph{Source and licence.} Minimal $a$-numbers of Artin--Schreier covers of ordinary curves. Version arXiv:2604.20091v1 (CC BY 4.0), distributed under the Creative Commons Attribution 4.0 International licence, as recorded in the arXiv licence metadata for this exact version and shown on the abstract page https://arxiv.org/abs/2604.20091v1. Full rights notice: licenses/arxiv-2604.20091-CC-BY-4.0.txt.\par\smallskip

\noindent\textit{Editorial scope.} Counted unit: the Thm whose source label is \texttt{thm:det} in the article (source0.tex lines 1058--1079, source label \texttt{thm:det}), delivered together with the article's own complete original proof of it (source0.tex lines 1081--1149, one \texttt{proof} environment of 69 lines). No mathematics is written, repaired or re-derived: statement and proof are the article's own text line for line. Required context retained uncounted: prop:M-props (lines 935--945); prop:phi (lines 972--981). Every definition, display and label of those lines is retained unchanged. Omitted from this package and not redistributed: 1--934, 946--971, 982--1057, 1080--1080, 1150--1197. Nothing inside a retained excerpt is omitted or altered: every retained body is delivered exactly as the article's lines are written, so no omission receipt of any kind accompanies this package. Citations: the retained excerpts carry no citation command at all, so no bibliography entry of the article is transcribed and no reference is redistributed. Reference adaptation: none was needed and none was made; every reference inside the delivered unit resolves inside it, so no unresolved reference or citation is produced. Printed slips kept literal: no printed word, symbol, hypothesis or number of the counted statement or of its proof was altered. Layout and class: the article's own class is \texttt{amsart} with packages including amssymb, amscd, amsrefs, xy, fullpage, libertine, tikz, fontenc; the wrapper is the lane's pinned \texttt{amsart} editorial wrapper at 10pt on A4, which restates the theorem-like environments the retained text needs and adds the article's own macro definitions that the retained text needs (\texttt{\textbackslash Cs}, \texttt{\textbackslash Rs}, \texttt{\textbackslash and}), with extra wrapper packages (bm, bbm, dsfont, xspace, cancel, mathtools). The delivered numbering is the unit's own. Assets: no retained span contains a figure environment, a graphics-inclusion command, a PDF-page inclusion, a table or a TikZ picture; no image file is copied into this package, the package declares no asset dependency and no image file is redistributed with it. Licence: version arXiv:2604.20091v1 of this article is distributed under the Creative Commons Attribution 4.0 International licence, as recorded in the arXiv licence metadata for this exact version and shown on the abstract page https://arxiv.org/abs/2604.20091v1. A full rights notice is delivered at licenses/arxiv-2604.20091-CC-BY-4.0.txt. Characters: every retained excerpt is pure ASCII, so the delivered engine, XeLaTeX with its default Latin Modern text font, reports no missing character.
\begin{prop}\label{prop:M-props}\mbox{}
  \begin{enumerate}
  \item If $M_{i,i'}=0$ \textup{(}i.e., if $pi-i'>\lambda d$\textup{)}
    then there is an $i''\in C_J$ such that $i''>i'$ and
    $M_{i,i''}=b_{\lambda,\lambda d}=a_d^\lambda$.
    In other words, the zero entries in $M$ all occur to the left of and
    in the same row as an occurrence of $b_{\lambda,\lambda d}$.  In
    particular, no row of $M$ is zero.
  \item  For all $((i,\lambda),i')\in R_J\times C_J$, we have $pi-i'>(\lambda-1)d$.
  \end{enumerate}
\end{prop}

\begin{prop}\label{prop:phi}
  We have:
  \begin{enumerate}
  \item $\phi(i,i')\equiv pi-i'\pmod d$
  \item The values of $\phi$ for a fixed $i$ and varying $i'$ are
    distinct.
  \item The values of $\phi$ for a fixed $i'$ and varying $i$ are also
    distinct.
  \end{enumerate}
\end{prop}

\begin{thm}\label{thm:det}\mbox{}
  \begin{enumerate}
  \item If $(i,i')\in T_N$, then
\[LM(N_{i,i'})=LM(b_{\lambda,pi-i'})=a_{\phi(i,i')}a_d^{\lambda-1}.\]
\item $Pr(\sigma)\ne0$ if and only if $\Gamma(\sigma)\subseteq T_N$.
\item If $\Gamma(\sigma)\subseteq T_N$, then
  \[LM(Pr(\sigma))=a_0^{e_0}a_1^{e_1}\cdots
    a_{d-1}^{e_{d-1}}a_d^{e_d+\Lambda}\]
  where for $1\le\ell\le d$,
  \[e_\ell=\left|\left\{i\in R_N|\phi(i,\sigma(i))=\ell\right\}\right|,\]
  and
  \[\Lambda=\sum_{i\in R_N}(\lambda-1).\]
  \textup{(}Recall that each index $i$ in $R_N$ determines a unique
  $\lambda$.\textup{)}
\item There exists a bijection $\sigma_0:R_N\to C_N$ such that
  $\Gamma(\sigma_0)\subseteq T_N$ and such that if
  $\sigma\neq\sigma_0$ is another bijection $\sigma:R_N\to C_N$ with
  $\Gamma(\sigma)\subseteq T_N$, then
  $LM(Pr(\sigma_0))>LM(Pr(\sigma))$.
\item $\det(N)\in R$ is non-zero.
  \end{enumerate}
\end{thm}

\begin{proof}\mbox{}
  (1) Note that $b_{\lambda,pi-i'}$ is homogeneous of degree $\lambda$
  and has weight $pi-i'=(\lambda-1)d+\phi(i,i')$.  Clearly
  $a_{\phi(i,i')}a_d^{\lambda-1}$ is the largest monomial of degree
  $\lambda$ and weight $(\lambda-1)d+\phi(i,i')$, and it appears in
  $f^\lambda$ with coefficient $\lambda$.

  (2) The product $Pr(\sigma)$ is non-zero if and only if each of the factors is
  non-zero if and only if $\Gamma(\sigma)\subseteq T_N$.

  (3) If $\Gamma(\sigma)\subseteq T_N$, then
  \[ LM(Pr(\sigma))=\prod_{i\in R} LM(N_{i,\sigma(i)}),\]
  and the result then follows from part (1).

  (4) We construct $\sigma_0$ with a greedy algorithm with steps
  indexed by $d,d-1,\dots,1$.  (See below for an example.)  At step
  $d$, consider
\[\Rs_d:=\left\{i\in R_N\mid\text{there is an $i'\in C_N$ with
      $\phi(i,i')=d$}\right\}.\]
For each $i\in \Rs_d$, part (2) of Proposition~\ref{prop:phi}
guarantees that there is a unique $i'$ such that $\phi(i,i')=d$, and
we set $\sigma_0(i)=i'$. This defines a map $\sigma_0:\Rs_d\to C_N$.
Part (1) of Proposition~\ref{prop:phi} guarantees that $\sigma_0$ is
injective.  Let $\Cs_d=\sigma_0(\Rs_d)$.  (If $\Rs_d$ happens to be
empty, we set $\Cs_d=\emptyset$ and proceed to the next step of the
algorithm.)

Note that Proposition~\ref{prop:M-props} (1) says that any zero
entries of $N$ occur in the same rows as entries where $\phi(i,i')=d$,
so if $i\in R_N\setminus \Rs_d$ and $i'\in C_N\setminus \Cs_d$, then
$N_{i,i'}\neq0$, i.e., $(i,i')\in T_N$ and $\phi$ is defined at
$(i,i')$.

In step $d-1$ of the algorithm, let
\[\Rs_{d-1}:=\left\{i\in R_N\setminus \Rs_d\mid\text{there is an $i'\in
      C_N\setminus \Cs_d$ with $\phi(i,i')=d-1$}\right\}.\]
Again, for each $i\in \Rs_{d-1}$, there is a unique
$C_N\setminus \Cs_d$ such that $\phi(i,i')=d-1$, and we set
$\sigma_0(i)=i'$.  (If $\Rs_{d-1}=\emptyset$, we set
$\Cs_{d-1}=\emptyset$ and proceed to the next step.)  This defines
$\sigma_0$ on $\Rs_d\cup \Rs_{d-1}$ and this map is injective.  Define
$\Cs_{d-1}:=\sigma_0(\Rs_{d-1})$.

We now continue, decreasing the sought after value of $\phi$ by one at
each step.  After $d$ steps, we have defined disjoint unions
\[R_N=\Rs_d\cup\cdots\cup\Rs_1
  \and
  C_N=\Cs_d\cup\cdots\cup\Cs_1,\]
  and a bijection $\sigma_0:R_N\to C_N$ such that
  $\sigma_0(\Rs_\ell)=\Cs_\ell$ and such that for $i\in \Rs_\ell$, we
  have $\phi(i,\sigma_0(i))=\ell$.

  Now suppose $\sigma\neq\sigma_0$ is another bijection $R_N\to C_N$.
  Let $\ell$ be the largest integer in $\{1,\dots,d\}$ such that
  $\sigma |_{\Rs_\ell}\neq\sigma_{0}|_{\Rs_\ell}$. Using part (3) of
  Theorem~\ref{thm:det}, we see that the exponents of
  $a_{\ell+1},\dots,a_d$ in $LM(Pr(\sigma))$ and $LM(Pr(\sigma_0))$
  are equal, but the exponent of $a_\ell$ in $LM(Pr(\sigma_0))$ is
  strictly larger than the exponent of $a_\ell$ in $LM(Pr(\sigma))$.
  This shows that $LM(Pr(\sigma_0))>LM(Pr(\sigma))$, as desired.

  (5) Part (4) shows that $LM(Pr(\sigma_0))$ is strictly larger than
  $LM(Pr(\sigma))$ for any bijection $\sigma\neq\sigma_0$ such that
  $Pr(\sigma)\neq0$.  Thus $LM(Pr(\sigma_0))$ cannot be cancelled by
  any monomial in $Pr(\sigma)$ for any $\sigma\neq\sigma_0$, and we
  conclude that $LM(Pr(\sigma_0))$ appears in $\det(N)$ with a
  non-zero coefficient.  This shows that $\det(N)\neq0$ and completes
  the proof of the Theorem.
\end{proof}

\end{document}
